\chapter{Of Trade, Chits, and Commerce}
\label{art:commerce}

\articlenote{In which is set down the chit, the silo's currency of trade, the marketplace where citizens buy and sell goods beyond the ration, and the keeping of accounts that make such commerce orderly and just.}

\section{The Chit and the Earning of Chits}

A chit is a token of value, issued by the silo and verified by the Chit-Keepers under Section 3, that represents a citizen's claim upon the silo's goods and the skilled labor of other citizens. A citizen earns chits through labor—by the work done in a trade, by the goods made or services rendered, by the time spent in the silo's keeping or the Departments' employ.

\begin{clauses}
\item Every citizen engaged in labor under Article 13 shall receive chits as payment for that labor, at a rate set by the Heads of Departments in consultation with the Mayor. The rate may vary by trade—skilled labor receives more than unskilled, dangerous labor may receive more than safe—but no citizen's chits shall be withheld or docked as punishment for work performed, save as Article 16 or 17 provides.
\item A citizen retired from labor under Article 13 shall continue to receive a modest sustenance in chits, sufficient to purchase small goods and comforts beyond the ration, administered by Supply under Article 11 upon the Infirmary's certification of the citizen's continuing need under Article 19.
\item A citizen sentenced to the mines under Article 17 receives no chits during the sentence; upon release, chits resume at the rate set for the citizen's assigned trade.
\item Chits are personal property and may not be seized, stolen, or diverted, save by Judicial order for the payment of debts or damages under Article 17. A citizen robbed of chits may appeal to the Sheriff for recovery; a citizen unable to recover them may petition Judicial for recompense.
\end{clauses}

\section{The Marketplace and the Sale of Goods}

Beyond the guaranteed ration of Article 11, citizens may purchase additional goods, specialized items, and the skilled labor of other citizens using chits. The marketplace is the ordinary mechanism of this trade: citizens may sell goods they have made or grown, or labor they have performed, to other citizens, with chits passing between them.

\begin{clauses}
\item A citizen may buy and sell goods with other citizens without restriction, save that no citizen may sell goods that belong to the silo or to a Department --- a Mechanical worker may not sell the silo's tools, nor a Supply clerk the silo's stores, and so for every Department alike.
\item A citizen engaged in a skilled trade—tailoring, cobbling, carpentry, or other craft—may sell the fruits of that trade to other citizens for chits, in addition to any goods the citizen's Department provides.
\item The silo maintains markets—designated areas on floors throughout the silo—where citizens may gather to buy and sell. The Sheriff under Article 7 keeps order in the markets and ensures no violence or theft occurs there.
\item A Market Clerk, raised up by short shadowing under Article 13 and appointed by the Mayor or the relevant Department Head, oversees a market and may set aside stalls for regular vendors or open vendors, may establish basic rules of conduct against spoiled goods and fraudulent claims, and may report disputes to Judicial or the Sheriff as required.
\end{clauses}

\section{The Chit-Keepers and the Regulation of Currency}

Chits are issued by the silo through the Chit-Keepers, who maintain the accounts of every citizen, verify the authenticity of chits in circulation, prevent counterfeiting, and ensure that the silo's currency holds its value and remains trustworthy.

\begin{clauses}
\item The Chit-Keepers are appointed by the Mayor and overseen by the Head of Supply under Article 11. The Chit-Keepers keep two rolls: one master roll of all chits issued, and one ledger of every citizen's balance, debits, and credits.
\item A citizen may, at any time, ask a Chit-Keeper to audit the citizen's balance or the authenticity of chits the citizen holds. The Chit-Keeper shall verify the chits and compare them to the citizen's ledger entry; any discrepancy is reported to Judicial.
\item Counterfeiting a chit, or possessing a counterfeit chit knowing it to be false, is among the gravest offenses under Article 17, as it strikes at the silo's foundation of trust in its own currency.
\item The Chit-Keepers may, at the Mayor's direction, revalue chits in the event of severe shortage or glut --- as, for example, where the silo's stores are depleted and chits have grown too numerous, the Mayor may order a revaluation downward. Such a revaluation is announced to the Assembly under Article 20 and entered in the Judicial archive.
\end{clauses}

\section{Prices, Fairness, and the Regulation of Commerce}

The silo does not fix the price of goods sold between citizens; a tailor and a customer may negotiate what chits a shirt is worth. But the silo does intervene where prices become extortionate or where a citizen attempts to monopolize a necessary good.

\begin{clauses}
\item A citizen may charge what price they wish for goods or labor offered to another citizen, and the other citizen may refuse to pay. No one is compelled to buy or sell at any given price.
\item If a citizen accuses another of charging an unjust price—unreasonably high compared to the goods' actual value or necessity—either party may appeal to the Market Clerk in the first instance; if either party remains unsatisfied, the matter is carried to Judicial, which alone may rule that the price is unjust and order a transaction reversed or the chits returned.
\item A citizen may not charge a price so high that other citizens are priced out of necessary goods --- as, for example, a water-seal seller charging so much that ordinary citizens cannot afford to seal their household water-stores. In such case the Mayor may, under the emergency authority of Article 22, set a maximum price for goods deemed essential; this is a measure of general policy reserved for genuine emergency, and is no substitute for Judicial's ordinary adjudication of a single unjust transaction under the clause above.
\item No citizen may prevent another citizen from selling goods or labor in the marketplace, save where Judicial finds the person selling is engaged in fraud, theft, or the sale of forbidden items under Article 16.
\end{clauses}

\section{Debt, Credit, and the Settling of Accounts}

A citizen may borrow chits from another citizen, creating a debt that must be repaid. The silo does not compel repayment of such debts, but the Chit-Keepers keep records of them, and Judicial may enforce them if a debtor refuses to pay and the creditor brings suit.

\begin{clauses}
\item A debt is a promise of repayment made between two citizens, recorded either by their own sworn statement to the Chit-Keepers or by a written note signed by both parties. A debt so recorded is legally binding.
\item A debtor who is unable to repay may petition Judicial for a reduction of the debt or an extended payment schedule; Judicial may grant such relief if the debtor's circumstances are genuine and the creditor is not unduly harmed.
\item A creditor who forgives a debt shall notify the Chit-Keepers, and the debt is removed from the ledger. A creditor who refuses to record the forgiveness, continuing to claim a debt the citizen has actually repaid, is answerable under Article 17 for fraud.
\item Enslavement to debt—a citizen being held to infinite servitude to pay a debt—is forbidden. No debt shall extend beyond the debtor's death, nor shall a debt be transferred to a debtor's children or heirs.
\end{clauses}

\section{The Goods and Services of the Silo}

The silo's skilled workers—tailors, cobblers, carpenters, rope-makers, and others—may be engaged for work by citizens willing to pay chits. The silo does not compel such work; the craftspeople may accept or refuse commissions as they choose.

\begin{clauses}
\item A skilled worker may, in addition to any work performed for their Department, undertake private commissions from citizens willing to pay chits. The worker and the citizen shall agree on a price and timeline before the work begins.
\item A citizen commissioned to make a good may refuse the commission without penalty, and a citizen commissioning a good may refuse to accept it or pay if the work is not as promised; disputes over such commissions are heard by the Market Clerk in the first instance, and, if either party remains unsatisfied, by Judicial.
\item Certain goods and services are restricted to licensed or skilled workers only --- as, for example, only a physician may sell medicines under Article 19, and only an electrician may repair certain technical goods under Article 8's licensing. A citizen advertising such goods or services without proper license is answerable under Article 17.
\end{clauses}

\section{The Protection of the Marketplace and the Silo's Trust}

The marketplace depends on trust: that chits are real, that sellers deliver what they promise, that buyers pay fairly, and that the Sheriff and the Chit-Keepers enforce these expectations. The protection of this trust is a standing charge of the Sheriff's office and the Chit-Keepers themselves.

\begin{clauses}
\item The Sheriff patrols the markets and nearby areas where citizens gather to trade, and may apprehend any citizen observed committing theft, fraud, or violence in the course of trade.
\item A citizen who is defrauded in a trade—sold counterfeit goods, shorted on chits, or given lesser work than promised—may report the matter to the Sheriff or the Market Clerk; that official shall investigate and, if fraud is found, refer the fraudster to Judicial.
\item The Chit-Keepers maintain security over the silo's stores of blank chits and the apparatus by which new chits are issued; any citizen who steals or tampers with blank chits, or who aids in their theft or counterfeiting, is answerable under Article 17 among the gravest offenses.
\end{clauses}
